\subsection{Factoring by Grouping} 
Sometimes, when we multiply two binomials, we get a polynomial in four terms. To factor such a polynomial, we want to get the two binomials back.
\begin{Example}
Multiply $(2x-3)(3y+1)$ by applying the distributive property twice.
\begin{lpic}{}{2}
\regtextleft{0.5}{-1}{$(2x-3)(3y+4)={}$};
\end{lpic}
\noindent Note that in this example, there are no \makeblank to combine.
\end{Example}
\noindent If we want to get back to the binomial factors, we can apply this process backward!
\begin{Example}
Factor $6xy-9y+8x-12$ as the product of two binomials.
\why[8]{5}
\end{Example}
\begin{procedure}{Factor by Grouping}
To factor a polynomial with four terms by grouping:
\begin{enumerate*}
\item Use the associative property of addition to ``group'' the terms into two pairs.
\item For each pair of terms, factor out the GCF.
\item Note that the remaining binomial factors for each pair are the same.
\item Factor out the remaining binomial factors to write the polynomial as a product of two binomials.
\item To check, multiply the two binomials, and confirm you get the original polynomial back.
\end{enumerate*}
\end{procedure}
Note that it is important that, when you create the groups, the operation combining them is \emph{addition}.
\begin{Example}
Factor $2y^3-3y^2-10y+15$ by grouping.
\begin{lpic}{}{4}
\regtextleft{0.5}{-1}{$2y^2-3y^2-10y+15={}$};
\end{lpic}
\end{Example}